\section{屋内で人の位置推定を行う手法に関する研究}
人の位置情報は様々な応用が期待されており，位置推定技術に対する需要は高い．
特に，人が日常生活や業務において屋内で過ごす時間は長く，屋内環境における位置情報の取得は重要な課題となっている．
近年，GPS を搭載したスマートフォンを利用した屋外位置推定が普及しているが，GPS は建物による遮蔽や反射の影響を受けやすく屋内環境において安定した位置の推定は困難である．
そのため，屋外位置推定における GPS に相当する，環境に依存せず広く適用可能である標準的な屋内位置推定手法が求められるが，未だ確立されていない．
こうした背景により，対象環境や利用目的に応じて様々なセンサや無線技術を用いた屋内位置推定手法が提案・研究されている．


屋内位置推定手法は，環境から直接位置を推定する絶対測位と，人の動作情報を用いて初期位置からの移動を推定する相対測位に分類され，それぞれに異なる利点と課題が存在する．
絶対測位の代表的な手法として，電波フィンガープリントを用いた位置推定や，複数の基準点からの距離情報を用いる三点測位方式が挙げられる．
電波フィンガープリントでは，事前に測位対象空間における Wi-Fi や BLE などの電波強度を含めた環境情報を収集し位置と対応させ，現在の測定値から位置を直接推定する\cite{fingerprint:WiFi}\cite{fingerprint:BLE}．
三点測位方式は位置が既知である複数の基地局やビーコンからの信号強度(RSSI)や到達時間をもとに距離を推定し，それらの交点として位置を算出する手法である\cite{threeposition:WiFi}\cite{threeposition:BLE}．
これらの絶対測位手法は，初期位置に依存せず位置推定が可能である反面，インフラの設置や環境変化への対応といった導入・運用コストが課題となる．
一方，相対測位の代表例として，加速度センサやジャイロセンサを用いた歩行者自律航法（PDR: Pedestrian Dead Reckoning）がある．
PDRは，初期位置を基準として歩行動作から移動量や方向を推定し，それらを順次足し合わせ現在位置を求める手法\cite{PDR:position}\cite{PDR:direction}である．
測位に特別なインフラを必要とせずスマートフォンなどの一般的な端末で利用可能である一方，推定誤差が時間とともに蓄積しやすいという課題がある．


屋内位置推定手法は測位精度の向上を図るほどインフラ設置や事前準備，運用管理にかかるコストが増大する傾向があり，精度とコストの間にはトレードオフが存在する．
しかし，屋内位置推定においては，利用目的や求められる測位精度に応じて，必ずしも高精度な位置推定が必要とは限らない．
例えば，在室管理や滞在状況の把握といった用途では位置を細かく推定する必要はなく，人物がどのエリアや部屋に存在するかを判別できれば十分な場合も多い．
このような目的に対して，BLE 信号強度に基づく近接判定（proximity）方式が提案，評価されている\cite{BLE:Proximity}．
位置精度は限定的であるものの，目的を限定する場合こうした導入や運用の負担が比較的小さい手法が選択される傾向にある．
他にも，単一の位置推定手法では精度や安定性に限界があるため，複数の手法を組み合わせて互いの欠点を補完するアプローチも多く用いられている．
例えば，相対測位であるPDRの誤差蓄積を抑制する手法として，外部情報を用いて推定位置を制約・補正するアプローチが用いられる場合が多い．
建物の構造を用いて物理的に不可能な経路とならないよう補正を行うマップマッチング\cite{mapmatching}や，既知の地点をランドマークとしてその地点で位置を再初期化するランドマーク測位\cite{landmark}が提案されている．
ただし，これらは事前準備や環境依存性といった新たなコストが発生する場合がある．